% ===================================================================
%  图表第 8 号 — 收获。— 打猎、摘葡萄、钓鱼。
%  Traduction 2026 d'apres texto/io, controlee sur texto/fr, texto/en
%  et texto/es. Consigne : voir texto/zh/10-tubiao-01.tex.
%
%  L'appel de note est dans le TITRE de la scene — « 收获 (*) » — et
%  non dans un alinea : c'est ainsi que l'ido le compose.
% ===================================================================

%%K t08-noto-1 noto
\VUnotes{143.2mm}{%
\textit{(*) 因为这个词不够确切：是收亚麻？摘葡萄？还是收苹果？也许用}
\VUgras{mesono} \textit{更好，见} Mondo \textit{第十四卷第 242 页。不过
这个新词根至今还没有得到学会采用，正等着按伊多语运动的通常规矩讨论。}%
}

%%K t08-apar-1 apar
\VUsaut{5.0mm}
\VUtitre{15.0pt}{60}{图表第 8 号}
\VUsaut{3.0mm}
\VUcentre{13.2pt}{90}{\textit{第一场。}}
\VUsaut{3.0mm}
\VUpk{10.2pt}{40}{收获 (*)。}
\VUsaut{4.0mm}
\VUpk{10.2pt}{40}{田野的面貌。}
\VUsaut{3.0mm}

%%K t08-c1-01-1 p
1. --- \VUgras{夏天}一到，田野又换了面貌。托付给垄沟的种子发了芽；地里长出
一株小小的绿苗，结了穗。谷粒长成，接着熟了，如今只剩下收割了。

%%K t08-c1-02-1 p
2. --- \VUgras{野花}开了。\VUgras{矢车菊} \textsuperscript{(1)}、\VUgras{罂粟}
\textsuperscript{(2)} 和\VUgras{毛地黄} \textsuperscript{(3)} 用它们新鲜的\VUgras{花}
给麦田那一色的底子添上活泼的对照。

%%K t08-c1-03-1 p
3. --- 果树长出了叶子；果子熟了。小小的\VUgras{樱桃树} \textsuperscript{(4)}
压满了好\VUgras{樱桃} \textsuperscript{(5)}，孩子们摘着，吃得津津有味。

%%K t08-c1-04-1 p
4. --- 羊群如今可以整天在外头了。

%%K t08-c1-04-2 p
\VUgras{羔羊} \textsuperscript{(6)} 长大了，成了毛厚厚的好\VUgras{母绵羊}
\textsuperscript{(7)} 和好\VUgras{公绵羊} \textsuperscript{(8)}。它们安安静静地在
\VUgras{牧场} \textsuperscript{(9)} 上啃\VUgras{草}，草里点缀着\VUgras{雏菊}。

%%K t08-c1-04-3 p
不久就要给这些牲口剪\VUgras{羊毛}，我们衣服的呢子就是用它做的。

%%K t08-tit-2 sub
\VUsaut{5.0mm}
\VUpk{10.2pt}{40}{放牛人。}
\VUsaut{3.4mm}

%%K t08-c1-05-1 p
5. --- \VUgras{放牛人} \textsuperscript{(10)} 好像不大理会它们。他心里只有天边
过去的那场雷雨；而那位\VUgras{牧羊人} \textsuperscript{(13)}，正把
\VUgras{水壶} \textsuperscript{(14)} 举到嘴边，看着更是无忧无虑。

%%K t08-c1-06-1 p
6. --- 天热得闷人；\VUgras{狗} \textsuperscript{(15)} 也来凉快凉快；它舔着
\VUgras{小溪} \textsuperscript{(16)} 里的清水，惊了一只本来蹲在岸边草里的可怜
小\VUgras{青蛙} \textsuperscript{(17)}。它跳进开着\VUgras{睡莲} \textsuperscript{(18)}
的清水里。

%%K t08-c1-07-1 p
7. --- 一只\VUgras{鹡鸰} \textsuperscript{(19)} 在小\VUgras{泉} \textsuperscript{(20)}
旁边洗澡，那泉从两块石头中间涌出来，钻到\VUgras{荆棘丛} \textsuperscript{(21)}
和\VUgras{山楂丛} \textsuperscript{(22)} 底下去了。

%%K t08-tit-3 sub
\VUsaut{5.0mm}
\VUpk{10.2pt}{40}{收割。}
\VUsaut{3.4mm}

%%K t08-c1-08-1 p
8. --- 对乡下孩子来说，割麦的时节很好玩。他们快活地\VUgras{翻筋斗}
\textsuperscript{(24)}，就在\VUgras{草堆} \textsuperscript{(23)} 上；他们帮着
\VUgras{割麦人} \textsuperscript{(26)}，这些人割着\VUgras{麦田} \textsuperscript{(27)}，
用的是\VUgras{镰刀} \textsuperscript{(25)}，在\VUgras{磨石} \textsuperscript{(30)}
上磨得很快；孩子们就把麦子拢起来，捆成\VUgras{麦捆} \textsuperscript{(31)} 或\VUgras{小把}
\textsuperscript{(32)}。

%%K t08-c1-09-1 p
9. --- 有时大些的孩子获准用\VUgras{镰刀} \textsuperscript{(12)} 割那些
\VUgras{金黄麦秆} \textsuperscript{(28)}，秆上托着饱满的\VUgras{麦穗}
\textsuperscript{(29)}；小些的孩子一边看着\VUgras{牲口} \textsuperscript{(33)}
在\VUgras{草地} \textsuperscript{(34)} 上、在大\VUgras{椴树} \textsuperscript{(35)}
荫下吃草，一边用他们的\VUgras{网} \textsuperscript{(36)} 捕漂亮的\VUgras{蝴蝶}。

%%K t08-c1-09-2 p
有几个甚至爬上\VUgras{大车} \textsuperscript{(37)}，去摆那些用\VUgras{草叉}
\textsuperscript{(38)} 递上来的麦捆。

%%K t08-c1-10-1 p
10. --- 整片\VUgras{平原} \textsuperscript{(39)} 上，\VUgras{云雀} \textsuperscript{(41)}
飞起来的地方，一派忙碌。人人都在忙收割。可是穷人还用着老家伙——
\VUgras{镰刀、钐刀}和\VUgras{连枷} \textsuperscript{(40)}——有钱的地主却割他们的麦子或
\VUgras{黑麦} \textsuperscript{(43)}，用的是\VUgras{收割机} \textsuperscript{(42)}。
然后也不用把麦捆运进谷仓，就地垛起大\VUgras{麦垛} \textsuperscript{(44)}。

%%K t08-c1-11-1 p
11. --- 接着开来一台\VUgras{脱粒机} \textsuperscript{(47)}，由一台
\VUgras{移动锅炉} \textsuperscript{(45)} 通过\VUgras{传动皮带} \textsuperscript{(46)} 带动，谷粒
就这样不离开播种的地就跟\VUgras{麦秆}分开了。

%%K t08-tit-4 sub
\VUsaut{5.0mm}
\VUpk{10.2pt}{40}{雷雨。}
\VUsaut{3.4mm}

%%K t08-c1-12-1 p
12. --- 收割时节常有猛烈的\VUgras{雷雨} \textsuperscript{(53)}。先是远处传来
\VUgras{雷}声；一阵强劲的\VUgras{西风} \textsuperscript{(54)} 刮起来，把
\VUgras{树林} \textsuperscript{(49)} 里最大的树吹得呻吟。黑\VUgras{云}
\textsuperscript{(56)} 攻占了天空；\VUgras{闪电} \textsuperscript{(55)} 刺眼的曲折
划破了它们。\VUgras{雨}下起来，有时还下\VUgras{冰雹}，把待割的庄稼打倒。

%%K t08-c1-13-1 p
13. --- 雷常常把房子引着。譬如去年，\VUgras{风磨} \textsuperscript{(50)} 的顶就
烧成了灰，两片\VUgras{磨翼} \textsuperscript{(51)} 也砸坏了。

%%K t08-c1-14-1 p
14. --- 过几个钟头，天又静了。天边现出一道灿烂的\VUgras{虹} \textsuperscript{(52)}，
万物又接上一时中断的生活。躲在\VUgras{窝棚} \textsuperscript{(48)} 里的乡下人
回去干活，勇敢地动手补救刚受的损失。

%%K t08-tit-5 sub
\VUsaut{6.0mm}
\VUcentre{13.2pt}{90}{\textit{第二场。}}
\VUsaut{3.6mm}

%%K t08-tit-6 sub
\VUpk{10.2pt}{40}{打猎。--- 摘葡萄。}
\VUsaut{1.4mm}
\VUpk{10.2pt}{40}{钓鱼。}
\VUsaut{4.0mm}

%%K t08-c2-01-1 p
1. --- 到\VUgras{八月}底或者\VUgras{九月}中，各地不同，狩猎期就开了。太阳
头几道光刚把\VUgras{蜥蜴} \textsuperscript{(2)} 引出洞来，\VUgras{猎人}
\textsuperscript{(5)} 就带着他的\VUgras{狗} \textsuperscript{(4)} 出门了。

%%K t08-c2-02-1 p
2. --- 他穿上厚呢子的结实\VUgras{猎装} \textsuperscript{(9)} 和皮\VUgras{绑腿}，
好走那藏着\VUgras{蟾蜍} \textsuperscript{(1)} 的湿草地；他拿了\VUgras{猎枪}
\textsuperscript{(6)} 和\VUgras{猎袋} \textsuperscript{(10)}，就上路了。

%%K t08-c2-03-1 p
3. --- 打猎的兴头把他引到一片葡萄园当中，那里正忙着摘葡萄。葡萄农的孩子
平常在路边看羊，今天却由着羊爱去哪儿吃就去哪儿；他们把一只老\VUgras{公山羊}
\textsuperscript{(3)} 拴在桩上——那家伙一得自由准会跑掉——自己也来凑个热闹。
一个从\VUgras{葡萄筐} \textsuperscript{(12)} 里拿了一串葡萄，挤出\VUgras{汁}
\textsuperscript{(14)} 到一只\VUgras{杯} \textsuperscript{(13)} 里，做成葡萄浆（没有
发酵的酒），觉得很有趣。另一个把刚编好的\VUgras{叶冠} \textsuperscript{(15)} 戴
在头上。他们旁边，大胆的\VUgras{麻雀} \textsuperscript{(16)} 一直飞到榨床跟前来
偷葡萄。

%%K t08-c2-04-1 p
4. --- 孩子玩着，大人干活。\VUgras{摘葡萄的人} \textsuperscript{(18)} 用
\VUgras{背篓} \textsuperscript{(17)} 把葡萄背到地窖去；一个\VUgras{箍桶匠}
\textsuperscript{(19)} 修\VUgras{酒桶} \textsuperscript{(20)}，查看\VUgras{桶板}
\textsuperscript{(22)} 和\VUgras{桶底} \textsuperscript{(21)} 是不是完好，换掉断了的
\VUgras{桶箍} \textsuperscript{(23)}。那些大\VUgras{缸} \textsuperscript{(25)} 也是他做的，
\VUgras{葡萄} \textsuperscript{(26)} 就倒在里面发酵。

%%K t08-c2-05-1 p
5. --- 发酵完了，就把酒放出来，剩下的渣放上\VUgras{榨床} \textsuperscript{(28)}；
一点一点拧紧那根大\VUgras{螺杆} \textsuperscript{(29)}，把汁挤出来，流进一只
\VUgras{槽} \textsuperscript{(32)} 里，又得了些\VUgras{酒} \textsuperscript{(31)}。

%%K t08-tit-8 sub
\VUsaut{6.0mm}
\VUpk{10.2pt}{40}{啤酒和苹果酒。}
\VUsaut{3.4mm}

%%K t08-c2-06-1 p
6. --- \VUgras{酒窖} \textsuperscript{(27)} 的墙上爬着一枝\VUgras{啤酒花}
\textsuperscript{(30)}。这里它只是观赏的，可是在有些国家，葡萄很少，就种啤酒花
来酿\VUgras{啤酒}。

%%K t08-c2-07-1 p
7. --- 小小的\VUgras{苹果树} \textsuperscript{(33)} 压满了苹果，熟了就能酿出好
\VUgras{苹果酒}。\VUgras{笼子} \textsuperscript{(34)} 里的\VUgras{金丝雀}
\textsuperscript{(35)} 和\VUgras{金翅雀} \textsuperscript{(36)} 眼馋地望着外面自由
自在蹦跳的麻雀。

%%K t08-c2-08-1 p
8. --- 一眼望不到头，山坡上排着一行行\VUgras{葡萄桩} \textsuperscript{(40)}，
撑着那些好\VUgras{葡萄藤} \textsuperscript{(39)}，藤上垂满了\VUgras{葡萄串}。

%%K t08-c2-08-2 p
一年到头\VUgras{葡萄农} \textsuperscript{(42)} 都在忙着照料他的\VUgras{葡萄园}
\textsuperscript{(38)}，如今好收成算是酬了他的辛苦。\VUgras{摘葡萄的女人}
\textsuperscript{(41)} 把\VUgras{骡子} \textsuperscript{(43)} 拉的车上的缸装满，人人
都高高兴兴。

%%K t08-tit-9 sub
\VUsaut{5.0mm}
\VUpk{10.2pt}{40}{钓鱼。}
\VUsaut{3.4mm}

%%K t08-c2-09-1 p
9. --- \VUgras{钓鱼的人} \textsuperscript{(44)} 也一样，天一亮就带着\VUgras{钓竿}
\textsuperscript{(45)} 来到水边坐下。

%%K t08-c2-09-2 p
\VUgras{鱼} \textsuperscript{(47)} 一见他们的\VUgras{钓线} \textsuperscript{(46)}
入水就上钩；他们刚来得及换上\VUgras{鱼饵}，线头上就又有一条在挣扎了。

%%K t08-c2-09-3 p
不过那个\VUgras{撒网人} \textsuperscript{(48)} 捕得还要多得多，他从
\VUgras{船} \textsuperscript{(49)} 上往\VUgras{河} \textsuperscript{(50)} 中撒网。

%%K t08-c2-10-1 p
10. --- 秋天是出游的好时候：步行、\VUgras{骑马} \textsuperscript{(52)}、骑
\VUgras{自行车} \textsuperscript{(53)}、坐\VUgras{汽车} \textsuperscript{(54)}。
\VUgras{路} \textsuperscript{(51)} 上干净，人们趁着最后几个好天出门，因为
\VUgras{燕子} \textsuperscript{(56)} 已经在屋顶上聚集，就要飞往暖和的地方了。
老\VUgras{核桃树} \textsuperscript{(57)} 的叶子黄了，\VUgras{核桃}落着，
成\VUgras{丛} \textsuperscript{(58)} 立在\VUgras{高地} \textsuperscript{(59)}
上的高\VUgras{树}也快要光秃了。

%%K t08-c2-11-1 p
11. --- \VUgras{太阳} \textsuperscript{(60)} 起得晚了，白天短了，\VUgras{早上}
已经觉得相当凉了，一群群\VUgras{丘鹬} \textsuperscript{(61)} 也已经离开我们这
不好客的地方了。
\VUsaut{6.0mm}
\VUfilet{22mm}
